\chapter{Corneal Transplant Surgery}
\index{Corneal Transplant Surgery}

\section*{Definition and Overview}
Corneal transplant surgery is the procedure in which damaged or diseased corneal tissue is replaced with healthy donor tissue. The cornea is the clear front window of the eye, and when it becomes scarred, cloudy, or misshapen, vision is lost. The surgery is performed by an ophthalmologist, usually under local anaesthetic with sedation or a general anaesthetic, and takes about one to two hours. Modern techniques allow some transplants to be performed through very small incisions, replacing only the diseased layer, which speeds recovery and reduces complications. This chapter describes what happens before, during, and after corneal transplant surgery.

\section*{Epidemiology}
Corneal transplantation is one of the most commonly performed tissue transplants in the world. In many countries and comparable countries, several thousand procedures are carried out each year. The most frequent indications are keratoconus, corneal scarring from infection or injury, and endothelial failure after cataract surgery. Demand for donor corneas is met through eye banks, and waiting times are usually short because corneal donation is well established. Outcomes are generally excellent, with clear grafts at one year in more than 90\% of cases for the most common indications.

\section*{Aetiology and Pathophysiology}
The surgical approach is determined by the layer of the cornea that is diseased.

\begin{itemize}
    \item \textbf{Full-thickness (penetrating) keratoplasty:} the entire central cornea is removed and replaced with donor tissue stitched into place
    \item \textbf{Deep anterior lamellar keratoplasty:} only the outer and middle layers are replaced, leaving the patient's own healthy endothelium
    \item \textbf{Endothelial keratoplasty:} the inner layer is replaced through a small side incision, without corneal stitches in most cases
    \item \textbf{Donor preparation:} the eye bank prepares the tissue to the correct thickness and size for the planned technique
    \item \textbf{Tissue matching:} blood group and immune matching are usually not required for the cornea, although screening for infection is routine
\end{itemize}

\section*{Clinical Features}
Before surgery, patients typically describe worsening vision that cannot be corrected with glasses.

\begin{itemize}
    \item Blurring or clouding of vision over months or years
    \item Distorted vision, glare, and halos around lights
    \item A history of keratoconus, corneal infection, injury, or previous eye surgery
    \item Eye pain or a sensation of pressure in conditions such as bullous keratopathy
    \item A visible haze or white area on the front of the eye
\end{itemize}

\section*{Differential Diagnosis}
The decision to operate is made only after confirming that the corneal problem is the main cause of reduced vision.

\begin{itemize}
    \item Cataract, where clouding of the lens causes similar symptoms
    \item Retinal disease, such as macular degeneration, reducing central vision
    \item Glaucoma, which can cause gradual visual field loss
    \item Refractive error treatable with glasses or contact lenses
    \item Corneal infection or inflammation that may respond to medical treatment alone
\end{itemize}

\section*{When to Seek Medical Care}
You will be reviewed by an ophthalmologist before surgery to confirm the diagnosis and plan the procedure. After surgery, urgent review is needed for any of the following.

\textbf{Contact your local emergency services for:}
\begin{itemize}
    \item Sudden severe eye pain or sudden loss of vision after surgery
    \item A red eye with discharge, which may indicate infection
    \item Injury to the operated eye
    \item Any sudden neurological symptoms such as weakness, difficulty speaking, or chest pain
\end{itemize}

\section*{Diagnosis and Investigations}
Pre-operative assessment confirms the diagnosis, establishes the extent of corneal disease, and ensures you are fit for surgery.

\begin{itemize}
    \item \textbf{History:} vision problems, previous eye surgery, infections, and general health
    \item \textbf{Slit-lamp examination:} detailed inspection of the cornea and its layers
    \item \textbf{Corneal topography and pachymetry:} mapping of the shape and thickness of the cornea
    \item \textbf{Endothelial cell assessment:} measurement of the health of the inner corneal layer
    \item \textbf{General health review:} blood pressure, blood tests, and medication review
    \item \textbf{Informed consent:} discussion of the risks, benefits, and expected recovery
\end{itemize}

\section*{Management}
The operation and recovery are managed by the ophthalmology team.

\begin{itemize}
    \item \textbf{Anaesthesia:} most commonly local anaesthetic with sedation, or general anaesthetic for children and complex cases
    \item \textbf{Surgical technique:} the diseased cornea is removed with a special circular blade or laser and replaced with donor tissue
    \item \textbf{Suturing:} full-thickness grafts are stitched into place with very fine stitches; endothelial grafts are usually held in place by an air bubble
    \item \textbf{Eye drops:} antibiotic and steroid drops are used for several months to prevent infection and rejection
    \item \textbf{Protective shield:} worn at night and when resting to prevent accidental rubbing
    \item \textbf{Follow-up:} frequent reviews in the first weeks, then less often over the following year
    \item \textbf{Stitch management:} sutures may be adjusted or removed as healing progresses
\end{itemize}

\section*{Complications}
\begin{itemize}
    \item Graft rejection, with redness, reduced vision, and sensitivity to light
    \item Infection inside the eye (endophthalmitis)
    \item Raised eye pressure (glaucoma)
    \item Bleeding or fluid leakage from the wound
    \item High astigmatism from uneven healing or tight stitches
    \item Failure of the graft to attach or clear
    \item Cataract formation
    \item Slow healing or wound gaping, more common with full-thickness grafts
\end{itemize}

\section*{Prognosis and Outcomes}
Most people leave hospital the same day or the next morning. Vision is often poor initially and improves gradually over several months as swelling settles and stitches are adjusted. Full visual recovery can take six to twelve months, and many patients need glasses or contact lenses afterwards. Graft survival is excellent: the majority of grafts remain clear at five years, and more than half are still functioning at ten years for favourable indications. Long-term monitoring remains important, as late rejection and other complications can still occur.

\section*{Prevention and Self-Care}
\begin{itemize}
    \item Use prescribed drops exactly as directed and do not stop them early
    \item Do not rub the eye; wear the protective shield at night
    \item Avoid swimming, heavy lifting, and strenuous exercise until cleared
    \item Attend every follow-up appointment, even if the eye feels well
    \item Report new pain, redness, or reduced vision without delay
    \item Protect the eye from injury during sport and gardening
    \item Keep underlying conditions such as glaucoma and herpes controlled
\end{itemize}

\section*{Sources and Further Reading}
The following reputable sources provide further information for healthcare students and patients seeking reliable, current advice about corneal transplant surgery and recovery.

\begin{itemize}
    \item NHS. \textit{Corneal transplant: the procedure and recovery.} https://www.nhs.uk/conditions/corneal-transplant/
    \item Mayo Clinic. \textit{Cornea transplant: surgery details and risks.} https://www.mayoclinic.org/tests-procedures/cornea-transplant/
    \item MSD Manual. \textit{Keratoplasty: surgical technique overview.} https://www.msdmanuals.com/
    \item MedlinePlus. \textit{Corneal transplant: what to expect.} https://medlineplus.gov/ency/article/003008.htm
    \item Royal Australian and New Zealand College of Ophthalmologists. \textit{Corneal transplant information.} https://ranzco.edu/
    \item Australian Government Organ and Tissue Authority. \textit{Corneal donation.} https://www.donatelife.gov.au/
\end{itemize}
